\documentclass[11pt,a4paper]{article}

% ==============================================================================
% CORE PACKAGES & TYPOGRAPHY
% ==============================================================================
\usepackage[utf8]{inputenc}
\usepackage[T1]{fontenc}
\usepackage{mathpazo}                 % Academic Palatino typography
\usepackage[scaled=0.92]{helvet}      % Clean sans-serif headings
\usepackage{inconsolata}              % Clean monospaced code font
\usepackage{microtype}                % Micro-typographic refinements
\usepackage[top=0.9in,bottom=0.9in,left=1in,right=1in,headheight=14pt]{geometry}
\usepackage{setspace}
\setstretch{1.12}                     % Slightly tighter line spacing

% ==============================================================================
% GRAPHICS, TABLES & TIKZ
% ==============================================================================
\usepackage{graphicx}
\usepackage{booktabs}
\usepackage{tabularx}
\usepackage{array}
\usepackage{amsmath,amssymb}
\usepackage{fancyhdr}
\usepackage{titlesec}
\usepackage{tikz}
\usetikzlibrary{shapes.geometric, arrows.meta, positioning, calc}

% ==============================================================================
% CODE LISTINGS STYLING (Subtle, Monochrome/Academic)
% ==============================================================================
\usepackage{listings}
\usepackage{xcolor}

\definecolor{codebg}{gray}{0.97}
\definecolor{codeframe}{gray}{0.80}
\definecolor{codekeyword}{rgb}{0.05, 0.15, 0.55}
\definecolor{codecomment}{rgb}{0.35, 0.40, 0.45}
\definecolor{codestring}{rgb}{0.12, 0.45, 0.15}
\definecolor{codenumbers}{gray}{0.50}

\lstset{
    language=Java,
    backgroundcolor=\color{codebg},
    frame=single,
    rulecolor=\color{codeframe},
    basicstyle=\small\ttfamily,
    keywordstyle=\color{codekeyword}\bfseries,
    commentstyle=\color{codecomment}\itshape,
    stringstyle=\color{codestring},
    numbers=left,
    numberstyle=\tiny\color{codenumbers},
    numbersep=8pt,
    breaklines=true,
    breakatwhitespace=true,
    tabsize=4,
    showstringspaces=false,
    captionpos=b,
    xleftmargin=15pt,
    framexleftmargin=15pt
}

% ==============================================================================
% HYPERREF SETUP (Subtle Academic Links)
% ==============================================================================
\usepackage{hyperref}
\hypersetup{
    colorlinks=true,
    linkcolor=black,
    citecolor=black,
    urlcolor=black,
    pdfauthor={Aarush Dubey and Sommayadeep},
    pdftitle={Car Rental Management System - Project Report (CSE 307)},
    pdfsubject={SRM University-AP CSE 307 Course Project}
}

% ==============================================================================
% HEADERS & FOOTERS
% ==============================================================================
\pagestyle{fancy}
\fancyhf{}
\fancyhead[L]{\small\nouppercase{\leftmark}}
\fancyhead[R]{\small\textit{CSE 307: Advanced Java Programming}}
\fancyfoot[C]{\thepage}
\renewcommand{\headrulewidth}{0.4pt}
\renewcommand{\footrulewidth}{0.0pt}

% Section title formatting
\titleformat{\section}{\large\bfseries}{\thesection}{1em}{}
\titleformat{\subsection}{\normalsize\bfseries}{\thesubsection}{1em}{}
\titleformat{\subsubsection}{\small\bfseries}{\thesubsubsection}{1em}{}

% Tighter spacing around sections
\titlespacing*{\section}{0pt}{1.2ex plus 0.5ex}{0.8ex plus 0.3ex}
\titlespacing*{\subsection}{0pt}{1.0ex plus 0.4ex}{0.6ex plus 0.2ex}
\titlespacing*{\subsubsection}{0pt}{0.8ex plus 0.3ex}{0.4ex plus 0.2ex}

% ==============================================================================
% DOCUMENT BEGIN
% ==============================================================================
\begin{document}

% ------------------------------------------------------------------------------
% FORMAL TITLE PAGE (Page 1)
% ------------------------------------------------------------------------------
\begin{titlepage}
    \centering
    \vspace*{0.5cm}
    
    {\Large\textbf{SRM UNIVERSITY-AP, ANDHRA PRADESH}}\\[0.25cm]
    {\large Department of Computer Science and Engineering}\\[0.15cm]
    {\small School of Engineering and Sciences}\\[1.2cm]
    
    \rule{\linewidth}{1.2pt}\\[0.4cm]
    {\LARGE\textbf{Design and Implementation of a Concurrent, Persistent Vehicle Fleet and Car Rental Management System}}\\[0.35cm]
    {\large\textit{Laboratory Course Project Report on Object-Oriented Architecture, Multithreaded Concurrency, JDBC Persistence, Generics, and Graphical User Interfaces}}\\[0.2cm]
    \rule{\linewidth}{1.2pt}\\[1.0cm]
    
    \textbf{Course Code \& Title:} Advanced Java Programming (\textbf{CSE 307})\\[1.5cm]
    
    \begin{minipage}{0.48\textwidth}
        \begin{flushleft}
            \textbf{Submitted by:}\\[0.2cm]
            \textbf{Aarush Dubey}\\
            Registration No: \texttt{AP24110010270}\\[0.2cm]
            \textbf{Sommayadeep}\\
            Registration No: \texttt{AP24110010181}\\[0.3cm]
            \textit{Bachelor of Technology}\\
            Department of Computer Science and Engineering\\
            SRM University-AP
        \end{flushleft}
    \end{minipage}
    \hfill
    \begin{minipage}{0.48\textwidth}
        \begin{flushright}
            \textbf{Submitted to:}\\[0.2cm]
            \textbf{Prof. Gokaramaiah Thota}\\
            Assistant Professor\\
            Department of Computer Science and Engineering\\
            School of Engineering and Sciences\\
            SRM University-AP
        \end{flushright}
    \end{minipage}
    
    \vfill
\end{titlepage}

% ------------------------------------------------------------------------------
% DECLARATION & CERTIFICATE (Page 2)
% ------------------------------------------------------------------------------
\newpage
\thispagestyle{plain}
\section*{Student Declaration}

We hereby declare that the laboratory project entitled \textbf{``Design and Implementation of a Concurrent, Persistent Vehicle Fleet and Car Rental Management System''} submitted in partial fulfillment of the requirements for the laboratory evaluation of \textbf{CSE 307: Advanced Java Programming} is an authentic record of original engineering work carried out by us under the academic supervision and guidance of \textbf{Prof. Gokaramaiah Thota}, Department of Computer Science and Engineering, School of Engineering and Sciences, SRM University-AP.

The design specifications, object-oriented domain models, multithreaded synchronization constructs, database integrations, test evaluations, and documentation contained in this report were developed and verified by us. No part of this report or code has been plagiarized or submitted previously for the award of any degree or diploma. All external tools, libraries, and references have been cited with proper academic attribution.

\vspace{2.2cm}

\noindent
\begin{minipage}{0.48\textwidth}
    \begin{flushleft}
        \rule{6cm}{0.4pt}\\
        \textbf{Aarush Dubey}\\
        Reg. No: \texttt{AP24110010270}\\
        Date: \today
    \end{flushleft}
\end{minipage}
\hfill
\begin{minipage}{0.48\textwidth}
    \begin{flushright}
        \rule{6cm}{0.4pt}\\
        \textbf{Sommayadeep}\\
        Reg. No: \texttt{AP24110010181}\\
        Date: \today
    \end{flushright}
\end{minipage}

\vspace{2.2cm}

\section*{Certificate of Evaluation}

This is to certify that this laboratory project report entitled \textbf{``Design and Implementation of a Concurrent, Persistent Vehicle Fleet and Car Rental Management System''} submitted by \textbf{Aarush Dubey} (Registration No: \texttt{AP24110010270}) and \textbf{Sommayadeep} (Registration No: \texttt{AP24110010181}) has been examined and evaluated by the department. The project satisfies all curricular expectations and assessment rubrics for \textbf{CSE 307: Advanced Java Programming}.

\vspace{2.0cm}

\noindent
\begin{minipage}{0.48\textwidth}
    \begin{flushleft}
        \rule{6.5cm}{0.4pt}\\
        \textbf{Prof. Gokaramaiah Thota}\\
        Course Instructor \& Project Supervisor\\
        Department of CSE, SRM University-AP
    \end{flushleft}
\end{minipage}
\hfill
\begin{minipage}{0.48\textwidth}
    \begin{flushright}
        \rule{6.5cm}{0.4pt}\\
        \textbf{Head of the Department / Examiner}\\
        Department of Computer Science and Engineering\\
        SRM University-AP
    \end{flushright}
\end{minipage}

% ------------------------------------------------------------------------------
% ABSTRACT & ACKNOWLEDGMENTS (Page 3)
% ------------------------------------------------------------------------------
\newpage
\thispagestyle{plain}
\section*{Abstract}

This laboratory project presents the design, implementation, and verification of a standalone \textbf{Vehicle Fleet and Car Rental Management System} in Java. Developed in accordance with the five curriculum units of \textbf{CSE 307: Advanced Java Programming}, the software integrates:
\begin{enumerate}
    \item \textbf{Object-Oriented Abstraction and Inheritance:} The \texttt{Rentable} interface, an abstract \texttt{Vehicle} base class, a concrete \texttt{RentalCar} subclass with static fleet tracking and immutable agency metadata, and a specialized \texttt{LuxurySedan} demonstrating multilevel inheritance.
    \item \textbf{Polymorphism and Exception Handling:} Dynamic method dispatch via V-Table resolution, static method overloading, and custom checked exceptions (\texttt{VehicleNotFoundException}) managed by \texttt{try-catch-finally} blocks.
    \item \textbf{Multithreading and Byte Streams:} Threads via \texttt{Thread} subclassing and \texttt{Runnable}, monitor synchronization (\texttt{wait()}/\texttt{notify()}), and binary streaming via \texttt{FileOutputStream}/\texttt{FileInputStream}.
    \item \textbf{JDBC Persistence and Collections:} Parameterized SQL via SQLite JDBC with \texttt{PreparedStatement}, synchronized with \texttt{ArrayList}, \texttt{HashSet}, and \texttt{HashMap} collections.
    \item \textbf{Generics and GUI:} Custom parameterized classes (\texttt{Pair<K,V>}), bounded wildcards (\texttt{? extends Vehicle}), and a Swing CRUD desktop GUI with headless detection.
\end{enumerate}
Testing confirms deterministic operation, zero data corruption under concurrency, and consistent behavior across desktop and headless environments.

\vspace{1.2cm}

\section*{Acknowledgments}

We express our gratitude to our course instructor and project supervisor, \textbf{Prof. Gokaramaiah Thota}, Assistant Professor, Department of Computer Science and Engineering, SRM University-AP. His guidance on JVM internals, concurrency semantics, and defensive software engineering laid the foundation for this implementation. We also thank the Department of CSE and School of Engineering and Sciences for providing the computing infrastructure essential for this project.

% ------------------------------------------------------------------------------
% TABLE OF CONTENTS (Page 4)
% ------------------------------------------------------------------------------
\newpage
\tableofcontents
\newpage

% ------------------------------------------------------------------------------
% RUBRIC MAPPING & EXECUTIVE OVERVIEW (Page 5)
% ------------------------------------------------------------------------------
\section*{Evaluation Rubric Cross-Reference Matrix}
\addcontentsline{toc}{section}{Evaluation Rubric Cross-Reference Matrix}

Table~\ref{tab:rubric_mapping} maps each criterion from the official department evaluation rubric to the corresponding sections and implementations in this report.

\begin{table}[h!]
\centering
\small
\begin{tabularx}{\textwidth}{p{3.8cm} p{6.2cm} p{3.8cm}}
\toprule
\textbf{Rubric Criteria} & \textbf{Key Implementations} & \textbf{Report Section} \\
\midrule
\textbf{1. Understanding of Java Program Structure} & Encapsulation, abstract base class, interface contract, inheritance hierarchy, static/final, constructor chaining. & Section~\ref{sec:structure} \\
\addlinespace
\textbf{2. Correct Use of Java Syntax} & Method overriding/overloading, custom checked exceptions, generics, collections, JDBC parameterized queries. & Section~\ref{sec:syntax} \\
\addlinespace
\textbf{3. Program Logic and Implementation} & CRUD workflow, tier pricing, multithreading, monitor synchronization, byte streams, Swing event model. & Section~\ref{sec:logic} \\
\addlinespace
\textbf{4. Source File Management} & Single-file compilation unit, vendored dependencies, schema bootstrap. & Section~\ref{sec:source} \\
\addlinespace
\textbf{5. Compilation and Execution} & Cross-platform classpath, shell/batch scripts, headless detection. & Section~\ref{sec:compilation} \\
\addlinespace
\textbf{6. Output and Theory} & Syllabus Units 1--5 mapping, console logs, database verification, test matrix. & Section~\ref{sec:output} \\
\bottomrule
\end{tabularx}
\caption{Cross-reference matrix linking evaluation rubrics to report sections.}
\label{tab:rubric_mapping}
\end{table}

\subsection*{Executive Architectural Overview}
The software follows a 3-tier decomposition:
\begin{itemize}
    \item \textbf{Presentation Tier:} Java Swing components (\texttt{JFrame}, \texttt{JTextField}, \texttt{JButton}) with lambda action listeners.
    \item \textbf{Application Tier:} Vehicle models, pricing logic, multithreaded dispatchers, and generic containers within \texttt{Main.java}.
    \item \textbf{Data Tier:} SQLite database (\texttt{rentals.db}) via JDBC combined with binary byte streams (\texttt{rental\_agreement.dat}).
\end{itemize}

% ==============================================================================
% SECTION 1: UNDERSTANDING OF JAVA PROGRAM STRUCTURE
% ==============================================================================
\newpage
\section{Understanding of Java Program Structure}
\label{sec:structure}

The foundational architecture is structured around core OOP paradigms. This section details interface abstraction, class hierarchies, constructor dispatch, and memory layout within the JVM.

\subsection{Object-Oriented Domain Architecture}
Vehicles share foundational properties (inventory IDs, manufacturer designations) but differ in pricing, capacities, and service options. The system models these through an extensible class hierarchy anchored by behavioral interfaces and abstract data models. Figure~\ref{fig:class_hierarchy} illustrates the inheritance structure.

\begin{figure}[h!]
\centering
\begin{tikzpicture}[
    node distance=1.4cm and 1.8cm,
    block/.style={rectangle, draw=black, thick, fill=gray!8, align=center, minimum height=1.0cm, text width=4.8cm},
    interface_block/.style={rectangle, draw=black, thick, dashed, fill=white, align=center, minimum height=0.9cm, text width=4.8cm},
    arrow/.style={-{Latex[scale=1.1]}, thick}
]
    \node[interface_block] (rentable) {
        \textbf{\guillemotleft interface\guillemotright}\\
        \textbf{Rentable}\\
        \small\texttt{+ displayRentalInfo(): void}
    };
    \node[block, below=of rentable] (vehicle) {
        \textbf{\textit{abstract Vehicle}}\\
        \footnotesize\texttt{\# vehicleId: int}\\
        \footnotesize\texttt{\# modelName: String}\\
        \footnotesize\texttt{+ \textit{getVehicleTier(): String}}
    };
    \node[block, below=of vehicle] (rentalcar) {
        \textbf{RentalCar}\\
        \footnotesize\texttt{+ static totalFleet: int}\\
        \footnotesize\texttt{+ static final AGENCY: String}\\
        \footnotesize\texttt{\# category, dailyRate}
    };
    \node[block, below=of rentalcar] (luxurysedan) {
        \textbf{LuxurySedan}\\
        \footnotesize\texttt{\# chauffeurIncluded: boolean}
    };
    \draw[arrow, dashed] (rentable) -- node[right, font=\footnotesize] {implements} (vehicle);
    \draw[arrow] (vehicle) -- node[right, font=\footnotesize] {extends} (rentalcar);
    \draw[arrow] (rentalcar) -- node[right, font=\footnotesize] {extends} (luxurysedan);
\end{tikzpicture}
\caption{Class and Interface Inheritance Architecture of the Car Rental System.}
\label{fig:class_hierarchy}
\end{figure}

\subsection{Interface Contract vs. Abstract State}
A key design decision is choosing between an \texttt{interface} and an \texttt{abstract class}. Both are used here strategically:

The \textbf{\texttt{Rentable} Interface} defines a purely behavioral contract --- if the enterprise later introduces non-vehicular rental assets (GPS navigators, child seats), those classes can implement \texttt{Rentable} without inheriting vehicular state:
\begin{lstlisting}[caption={The Rentable Interface Contract}]
interface Rentable {
    void displayRentalInfo();
}
\end{lstlisting}

The \textbf{\texttt{abstract Vehicle} Class} represents an incomplete entity. All vehicles share an inventory ID and model name, so these are defined with \texttt{protected} access in an abstract class. Direct instantiation via \texttt{new Vehicle(...)} is blocked since it lacks tariff data:
\begin{lstlisting}[caption={Abstract Vehicle Base Class}]
abstract class Vehicle implements Rentable {
    protected int vehicleId;
    protected String modelName;

    public Vehicle(int vehicleId, String modelName) {
        this.vehicleId = vehicleId;
        this.modelName = modelName;
    }
    public abstract String getVehicleTier();

    public void displayRentalInfo() {
        System.out.println("ID: " + vehicleId + ", Model: " + modelName
            + " [" + getVehicleTier() + "]");
    }
}
\end{lstlisting}

\subsection{Inheritance Hierarchy and Multilevel Specialization}
The domain implements a three-tier multilevel hierarchy:
\begin{enumerate}
    \item \textbf{Tier 1 --- \texttt{Vehicle}:} Core identity fields; enforces \texttt{Rentable} implementation.
    \item \textbf{Tier 2 --- \texttt{RentalCar}:} Adds commercial parameters (\texttt{category}, \texttt{dailyRate}), provides \texttt{getVehicleTier()} returning \texttt{"Standard Rental Car"}, and implements fuel efficiency logic.
    \item \textbf{Tier 3 --- \texttt{LuxurySedan}:} Adds a chauffeur flag (\texttt{boolean chauffeurIncluded}) and overrides \texttt{getVehicleTier()} to dynamically report driver inclusion --- exhibiting runtime polymorphism.
\end{enumerate}

\subsection{Static/Final Qualifiers and Constructor Chaining}
The project uses \texttt{static} and \texttt{final} to separate instance and class-level data:
\begin{itemize}
    \item \textbf{\texttt{public static int totalFleet}:} A shared counter in JVM Metaspace, incremented by every constructor invocation across all subclasses.
    \item \textbf{\texttt{public static final String AGENCY}:} Initialized to \texttt{"SRM Rentals"} and immutable after compilation, enabling constant folding and thread-safe reads.
\end{itemize}

Constructor chaining ensures no object exists in an incomplete state:
\begin{lstlisting}[caption={Constructor Chaining in RentalCar and LuxurySedan}]
class RentalCar extends Vehicle {
    public RentalCar(int vehicleId, String modelName,
                     String category, double dailyRate) {
        super(vehicleId, modelName);
        this.category = category;
        this.dailyRate = dailyRate;
        totalFleet++;
    }
    // Overloaded: delegates with default rate
    public RentalCar(int vehicleId, String modelName, String category) {
        this(vehicleId, modelName, category, 1500.0);
    }
}

class LuxurySedan extends RentalCar {
    public LuxurySedan(int id, String model, String cat,
                       double rate, boolean chauffeur) {
        super(id, model, cat, rate);
        this.chauffeurIncluded = chauffeur;
    }
}
\end{lstlisting}

When \texttt{new RentalCar(105, "Nexon", "SUV")} is called: (1)~the 3-arg constructor delegates to the 4-arg via \texttt{this(..., 1500.0)}, (2)~the 4-arg calls \texttt{super(vehicleId, modelName)}, (3)~the parent sets identity fields, (4)~control returns to set \texttt{category}, \texttt{dailyRate}, and increment \texttt{totalFleet}.

% ==============================================================================
% SECTION 2: CORRECT USE OF JAVA SYNTAX
% ==============================================================================
\newpage
\section{Correct Use of Java Syntax}
\label{sec:syntax}

This section evaluates the code against Java language specifications, type safety, exception handling, and defensive programming conventions.

\subsection{Polymorphism: Compile-Time vs. Runtime Dispatch}

\subsubsection{Dynamic Method Dispatch (Overriding)}
When a method is called on a base reference, the JVM resolves the actual implementation at runtime via the V-Table:
\begin{lstlisting}[caption={Dynamic Method Dispatch via Base Reference}]
Vehicle polyVehicle = new LuxurySedan(202, "BMW 5 Series",
                                       "Executive", 7200.0, true);
System.out.println("Tier: " + polyVehicle.getVehicleTier());
// Output: Executive Luxury Sedan (Chauffeur: Included)
\end{lstlisting}
The compiler validates that \texttt{Vehicle} declares \texttt{getVehicleTier()}. At runtime, the JVM inspects the object header, accesses the \texttt{LuxurySedan} V-Table, and executes the overridden method.

\subsubsection{Static Method Overloading}
Compile-time polymorphism is demonstrated through overloaded search routines:
\begin{lstlisting}[caption={Overloaded Search with Differentiated Signatures}]
public static RentalCar search(List<RentalCar> list, int vehicleId)
    throws VehicleNotFoundException {
    for (RentalCar c : list)
        if (c.vehicleId == vehicleId) return c;
    throw new VehicleNotFoundException("Vehicle ID " + vehicleId
        + " not found!");
}

public static RentalCar search(List<RentalCar> list, String model)
    throws VehicleNotFoundException {
    for (RentalCar c : list)
        if (c.modelName.equalsIgnoreCase(model)) return c;
    throw new VehicleNotFoundException("Model '" + model
        + "' not found!");
}
\end{lstlisting}
The compiler distinguishes \texttt{(List, int)} from \texttt{(List, String)} by parameter types, resolving the target method at compile time with zero runtime overhead.

\subsection{Exception Handling Architecture}

\subsubsection{Custom Checked Exception}
Checked exceptions inherit from \texttt{java.lang.Exception}, forcing compile-time verification:
\begin{lstlisting}[caption={Custom Checked Domain Exception}]
class VehicleNotFoundException extends Exception {
    public VehicleNotFoundException(String msg) {
        super(msg);
    }
}
\end{lstlisting}

\subsubsection{Guaranteed Cleanup with \texttt{try-catch-finally}}
All search transactions are protected by structured blocks:
\begin{lstlisting}[caption={Try-Catch-Finally Execution}]
try {
    List<RentalCar> temp = new ArrayList<>();
    temp.add(c1);
    search(temp, 999); // Non-existent ID
} catch (VehicleNotFoundException e) {
    System.out.println("Caught: " + e.getMessage());
} finally {
    System.out.println("Finally block executed.");
}
\end{lstlisting}

\subsection{Generics and Bounded Wildcards}

A generic \texttt{Pair<K, V>} bundles paired attributes without single-purpose classes:
\begin{lstlisting}[caption={Generic Pair and Bounded Wildcard}]
class Pair<K, V> {
    public K key;
    public V value;
    public Pair(K k, V v) { this.key = k; this.value = v; }
}

// Upper-bounded wildcard: accepts List<RentalCar>, List<LuxurySedan>
public static void inspectFleet(List<? extends Vehicle> list) {
    for (Vehicle v : list)
        System.out.println("Inspected: " + v.modelName
            + " (" + v.getVehicleTier() + ")");
}
\end{lstlisting}
Following the PECS (Producer Extends, Consumer Super) rule, \texttt{List<? extends Vehicle>} acts as a producer of \texttt{Vehicle} references, safely accepting any subtype list.

\subsection{Java Collections Framework}
The system uses three collections based on their algorithmic properties (Table~\ref{tab:collections_comparison}):

\begin{table}[h!]
\centering
\small
\begin{tabularx}{\textwidth}{l l l p{5.0cm}}
\toprule
\textbf{Interface} & \textbf{Class} & \textbf{Complexity} & \textbf{Role in System} \\
\midrule
\texttt{List<E>} & \texttt{ArrayList} & Access: $O(1)$ & Sequential fleet storage and ordered iteration. \\
\addlinespace
\texttt{Set<E>} & \texttt{HashSet} & Add/Find: $O(1)$ & Unique vehicle categories (SUV, Sedan). \\
\addlinespace
\texttt{Map<K,V>} & \texttt{HashMap} & Get/Put: $O(1)$ & Instant lookup by \texttt{vehicleId}. \\
\bottomrule
\end{tabularx}
\caption{Algorithmic Comparison of Java Collections Used.}
\label{tab:collections_comparison}
\end{table}

\subsection{Defensive JDBC with PreparedStatements}
Database operations use \texttt{PreparedStatement} to prevent SQL injection. All connections use \texttt{try-with-resources} for guaranteed closure:
\begin{lstlisting}[caption={Parameterized JDBC Insertion}]
public static void saveVehicle(RentalCar c) throws Exception {
    try (Connection conn = getDbConnection()) {
        String sql = "INSERT OR REPLACE INTO fleet VALUES (?,?,?,?)";
        try (PreparedStatement ps = conn.prepareStatement(sql)) {
            ps.setInt(1, c.vehicleId);
            ps.setString(2, c.modelName);
            ps.setString(3, c.category);
            ps.setDouble(4, c.dailyRate);
            ps.executeUpdate();
        }
    }
}
\end{lstlisting}

% ==============================================================================
% SECTION 3: PROGRAM LOGIC AND IMPLEMENTATION
% ==============================================================================
\newpage
\section{Program Logic and Implementation}
\label{sec:logic}

This section details the CRUD workflow, multithreaded concurrency, inter-thread synchronization, byte streaming, and GUI event dispatching.

\subsection{Dual-State CRUD Workflow}
The software synchronizes state changes across the SQLite database and in-memory collections. Table~\ref{tab:crud_logic} outlines the execution logic.

\begin{table}[h!]
\centering
\small
\begin{tabularx}{\textwidth}{l p{2.8cm} p{4.0cm} p{4.5cm}}
\toprule
\textbf{Operation} & \textbf{Trigger} & \textbf{SQL} & \textbf{Memory Action} \\
\midrule
\textbf{Create} & \texttt{btnAdd} click & \texttt{INSERT OR REPLACE} & Remove stale ID, insert new object, update Map/Set. \\
\addlinespace
\textbf{Retrieve} & \texttt{btnSearch} click & \texttt{SELECT WHERE id=?} & Fetch from list, auto-fill UI fields. \\
\addlinespace
\textbf{Update} & \texttt{btnUpdate} click & \texttt{UPDATE SET rate=?} & Mutate \texttt{dailyRate} in memory. \\
\addlinespace
\textbf{Delete} & \texttt{btnDelete} click & \texttt{DELETE WHERE id=?} & \texttt{removeIf()} on list, purge from Map. \\
\bottomrule
\end{tabularx}
\caption{Dual-State CRUD Execution Logic.}
\label{tab:crud_logic}
\end{table}

\subsection{Dynamic Tariff Classification}
Vehicle efficiency categories are computed dynamically from pricing:
\begin{lstlisting}[caption={Fuel Efficiency Decision Logic}]
public String getFuelEfficiencyClass() {
    if (dailyRate >= 3500) return "Premium Performance";
    else if (dailyRate >= 2000) return "Comfort Cruiser";
    else return "Economy Mileage";
}
\end{lstlisting}

\subsection{Multithreading and Monitor Synchronization}
Fleet management requires that vehicle dispatch cannot occur until safety inspection completes. This uses Java's monitor locks and the guarded suspension pattern:

\begin{lstlisting}[caption={Thread Synchronization with Guarded Suspension}]
class FleetDispatchSync {
    private boolean inspectionDone = false;

    public synchronized void completeInspection() {
        inspectionDone = true;
        notify(); // Wake waiting thread
    }

    public synchronized void dispatchVehicle()
        throws InterruptedException {
        while (!inspectionDone) wait(); // Guard against spurious wakeup
        System.out.println("Safety inspection verified! Keys released.");
    }
}
\end{lstlisting}

The \texttt{while} loop (not \texttt{if}) guards against spurious wakeups --- a JVM behavior where threads may wake without an explicit \texttt{notify()} signal. The two thread creation mechanisms are also demonstrated:
\begin{itemize}
    \item \textbf{Extending \texttt{Thread}:} \texttt{FleetMaintenanceThread extends Thread} with an overridden \texttt{run()} method.
    \item \textbf{Implementing \texttt{Runnable}:} \texttt{DailyBillingTask implements Runnable}, passed to \texttt{new Thread(runnable).start()}.
\end{itemize}

\subsection{Binary Byte Stream I/O}
Low-level file I/O is demonstrated using \texttt{FileOutputStream} and \texttt{FileInputStream} to write and read raw bytes to \texttt{rental\_agreement.dat}, independent of character encodings.

\subsection{Swing GUI and Delegation Event Model}
The interactive CRUD interface uses Java Swing components (\texttt{JFrame}, \texttt{JTextField}, \texttt{JButton}, \texttt{JLabel}) organized under \texttt{FlowLayout}. User events are dispatched via lambda-based \texttt{ActionListener} implementations following the Delegation Event Model. Each CRUD button triggers the corresponding database operation and synchronizes the in-memory collections accordingly.

% ==============================================================================
% SECTION 4: SOURCE FILE MANAGEMENT
% ==============================================================================
\newpage
\section{Source File Management}
\label{sec:source}

\subsection{Project Directory Structure}
\begin{lstlisting}[language={}, caption={Project Directory Layout}]
java_project/
|-- Main.java            # Consolidated compilation unit (410 lines)
|-- rentals.db           # SQLite embedded database
|-- rental_agreement.dat # Binary byte stream manifest
|-- lib/
|   `-- sqlite-jdbc.jar  # Vendored JDBC driver
|-- compile.sh / run.sh  # POSIX shell scripts
|-- compile.bat / run.bat # Windows batch scripts
|-- RUN_COMMANDS.txt     # CLI reference
`-- README.md            # Project documentation
\end{lstlisting}

\subsection{Single-Unit Modular Consolidation}
All auxiliary classes (\texttt{Rentable}, \texttt{Vehicle}, \texttt{RentalCar}, \texttt{LuxurySedan}, \texttt{VehicleNotFoundException}, \texttt{FleetMaintenanceThread}, \texttt{DailyBillingTask}, \texttt{FleetDispatchSync}, \texttt{Pair}) are maintained as package-private entities alongside the public \texttt{Main} class within \texttt{Main.java}. This satisfies the JLS requirement of exactly one public class per source file while enabling single-command compilation without build tool overhead.

\subsection{Vendored Dependency and Schema Bootstrap}
The SQLite JDBC driver is vendored in \texttt{lib/sqlite-jdbc.jar}, bundling native libraries for macOS, Linux, and Windows. This guarantees offline execution during laboratory evaluations. The database schema is auto-created on startup:
\begin{lstlisting}[caption={Automated Schema Bootstrapping}]
st.executeUpdate(
    "CREATE TABLE IF NOT EXISTS fleet (" +
    "id INT PRIMARY KEY, model TEXT, category TEXT, rate REAL)");
\end{lstlisting}

% ==============================================================================
% SECTION 5: COMPILATION AND EXECUTION
% ==============================================================================
\section{Compilation and Execution}
\label{sec:compilation}

\subsection{Cross-Platform Classpath Configuration}
The classpath separator differs across platforms:
\begin{lstlisting}[language=bash, caption={macOS/Linux Commands}]
javac -cp "lib/sqlite-jdbc.jar" Main.java
java -cp ".:lib/sqlite-jdbc.jar" Main
\end{lstlisting}
\begin{lstlisting}[language={}, caption={Windows Commands}]
javac -cp "lib\sqlite-jdbc.jar" Main.java
java -cp ".;lib\sqlite-jdbc.jar" Main
\end{lstlisting}

\subsection{Automated Scripts and Headless Detection}
Shell scripts (\texttt{compile.sh}, \texttt{run.sh}) and batch files (\texttt{compile.bat}, \texttt{run.bat}) provide single-click compilation and execution. The application also detects headless environments:
\begin{lstlisting}[caption={Headless Environment Detection}]
if (!GraphicsEnvironment.isHeadless()) {
    JFrame f = new JFrame("Car Rental System");
    // ... configure and display GUI ...
    f.setVisible(true);
} else {
    System.out.println("[Notice] Headless environment. "
        + "Terminal demo completed.");
}
\end{lstlisting}
This allows the full test suite to execute in SSH sessions, CI pipelines, and terminal-only lab environments without \texttt{HeadlessException} crashes.

% ==============================================================================
% SECTION 6: OUTPUT AND THEORY
% ==============================================================================
\newpage
\section{Output and Theory}
\label{sec:output}

This section ties empirical output logs and test matrices back to the theoretical foundations across all five syllabus units.

\subsection{Theoretical Foundations Across Syllabus Units}

\textbf{Unit 1 --- OOP, Encapsulation, Inheritance:} Bundles data and methods within classes using access modifiers. The \texttt{RentalCar} ``is-a'' \texttt{Vehicle} relationship promotes code reuse; abstract classes provide partial implementations while interfaces define pure behavioral contracts.

\textbf{Unit 2 --- Polymorphism and Exceptions:} Dynamic dispatch resolves methods at runtime via V-Table indexing. Custom checked exceptions enforce compile-time error handling for anticipated operational faults.

\textbf{Unit 3 --- Multithreading and Byte Streams:} Demonstrates thread lifecycle states (NEW, RUNNABLE, WAITING, TERMINATED). Every Java object has an intrinsic monitor lock; \texttt{synchronized} serializes critical regions while \texttt{wait()}/\texttt{notify()} coordinate thread communication. Byte streams provide encoding-independent I/O.

\textbf{Unit 4 --- JDBC and Collections:} \texttt{DriverManager} loads drivers, establishes connections, and compiles SQL via \texttt{PreparedStatement}. The collections hierarchy contrasts indexed access (\texttt{List}), uniqueness (\texttt{Set}), and key-value mapping (\texttt{Map}).

\textbf{Unit 5 --- Generics and Swing GUI:} Type parameters prevent runtime \texttt{ClassCastException} errors. Upper-bounded wildcards enable covariance. Swing events are processed via the Delegation Event Model on the Event Dispatch Thread.

\subsection{Console Output Logs}
\begin{lstlisting}[language={}, caption={Application Console Output}]
=== VEHICLE / CAR RENTAL SYSTEM (CSE 307) ===

--- [UNIT 1] CLASSES & INHERITANCE ---
Agency: SRM Rentals
ID: 101, Model: Hyundai Creta [Standard Rental Car]
Category: SUV, Daily Rate: Rs.2800.0, Class: Comfort Cruiser
ID: 201, Model: Mercedes C-Class [Executive Luxury Sedan
  (Chauffeur: Included)]
Total Fleet Count (static): 2

--- [UNIT 2] POLYMORPHISM & EXCEPTIONS ---
Dynamic dispatch vehicle tier: Executive Luxury Sedan
  (Chauffeur: Included)
Caught Custom Exception: Vehicle ID 999 not found in fleet!
Finally block executed.

--- [UNIT 3] THREADS & BYTE STREAMS ---
[Thread 1] Fleet maintenance check running at Priority: 10
[Thread 2] Rental revenue audit finished via Runnable.
[Sync] Safety inspection verified via wait/notify!
Byte Stream Read: 101,Hyundai Creta,SUV,Rate:2800

--- [UNIT 4] JDBC & COLLECTIONS ---
Unique Categories (HashSet): [SUV, Luxury Sedan]
Fleet Map Keys (HashMap): [101, 201]
Overloaded Search (ID 101): Hyundai Creta
Overloaded Search (Model 'Hyundai Creta'): ID 101

--- [UNIT 5] GENERICS & GUI ---
Generic Pair (Most Exclusive): BMW 5 Series (Rs.7200.0/day)
Inspected: Hyundai Creta (Standard Rental Car)
Inspected: Mercedes C-Class (Executive Luxury Sedan
  (Chauffeur: Included))
Full CRUD GUI window launched!
\end{lstlisting}

\subsection{Database Verification}
The persistence layer was verified via the SQLite command-line shell:
\begin{lstlisting}[language=sql, caption={SQLite Verification Queries}]
sqlite> .schema fleet
CREATE TABLE fleet (id INT PRIMARY KEY, model TEXT,
                    category TEXT, rate REAL);
sqlite> SELECT * FROM fleet;
101|Hyundai Creta|SUV|2800.0
201|Mercedes C-Class|Luxury Sedan|6500.0
\end{lstlisting}

\subsection{System Verification Test Matrix}
Table~\ref{tab:verification_matrix} summarizes the formal test scenarios evaluated.

\begin{table}[h!]
\centering
\small
\begin{tabularx}{\textwidth}{l p{2.5cm} p{3.5cm} p{4.0cm} c}
\toprule
\textbf{ID} & \textbf{Component} & \textbf{Input / Action} & \textbf{Expected Result} & \textbf{Status} \\
\midrule
TC-01 & Class Hierarchy & Instantiate \texttt{RentalCar} & Fields populated; fleet count increments & \textbf{PASS} \\
\addlinespace
TC-02 & Polymorphism & \texttt{getVehicleTier()} on base ref & Dispatches to \texttt{LuxurySedan} override & \textbf{PASS} \\
\addlinespace
TC-03 & Overloading & Search by ID and String & Same vehicle found via both signatures & \textbf{PASS} \\
\addlinespace
TC-04 & Exceptions & Query non-existent ID 999 & \texttt{VehicleNotFoundException} caught & \textbf{PASS} \\
\addlinespace
TC-05 & Concurrency & Dispatch before inspection & Thread blocks in \texttt{wait()}; proceeds after \texttt{notify()} & \textbf{PASS} \\
\addlinespace
TC-06 & Byte I/O & Write \& read to \texttt{.dat} file & Data reconstructed without corruption & \textbf{PASS} \\
\addlinespace
TC-07 & JDBC & Insert via GUI Add button & Row persisted in \texttt{rentals.db} & \textbf{PASS} \\
\addlinespace
TC-08 & Headless & Run without display server & Detects headless; no crash & \textbf{PASS} \\
\bottomrule
\end{tabularx}
\caption{System Verification Test Matrix.}
\label{tab:verification_matrix}
\end{table}

% ==============================================================================
% SECTION 7: VIVA-VOCE
% ==============================================================================
\newpage
\section{Viva-Voce Oral Examination Preparation}
\label{sec:viva}

Concise answers to fundamental Java questions related to this project.

\subsection*{Q1: Why use an abstract class for \texttt{Vehicle} instead of an interface?}
\textbf{A:} All vehicles share state attributes (\texttt{vehicleId}, \texttt{modelName}). An abstract class encapsulates these with \texttt{protected} access and shares constructor logic via \texttt{super()}, while still enforcing abstract method implementation in subclasses.

\subsection*{Q2: Overloading vs. overriding in terms of dispatch?}
\textbf{A:} Overloading is static (compile-time): the compiler binds the method based on parameter types. Overriding is dynamic (runtime): the JVM resolves via the V-Table based on the actual object instance.

\subsection*{Q3: Why a checked exception instead of RuntimeException?}
\textbf{A:} A missing vehicle is an anticipated operational condition (user typo, unavailable car), not a programming bug. Checked exceptions enforce that callers implement deterministic recovery logic.

\subsection*{Q4: Why must \texttt{wait()} be in a \texttt{while} loop?}
\textbf{A:} Spurious wakeups can occur in JVM implementations. A \texttt{while} loop re-evaluates the condition predicate upon waking; if still false, the thread safely re-enters the wait state.

\subsection*{Q5: How does \texttt{PreparedStatement} prevent SQL injection?}
\textbf{A:} The database compiles the query template before parameter insertion. Parameters are transmitted as typed data values, not executable SQL fragments, rendering injections inert.

\subsection*{Q6: Explain the PECS rule.}
\textbf{A:} Producer Extends, Consumer Super. \texttt{inspectFleet(List<? extends Vehicle>)} reads (produces) \texttt{Vehicle} references. The upper bound allows covariance: lists of \texttt{RentalCar} or \texttt{LuxurySedan} can be safely passed.

% ==============================================================================
% SECTION 8: CONCLUSION
% ==============================================================================
\section{Conclusion and Future Enhancements}
\label{sec:conclusion}

The Car Rental Management System successfully covers all five curriculum units of \textbf{CSE 307}. The codebase demonstrates clean OOP architecture via abstract classes and interface contracts, runtime and compile-time polymorphism, defensive exception handling, concurrent thread management with monitor synchronization, persistent storage through JDBC and byte streams, and an interactive Swing GUI with headless detection.

\subsection{Future Enhancements}
\begin{enumerate}
    \item \textbf{Connection Pooling:} Migrating to HikariCP for improved throughput under concurrent database access.
    \item \textbf{Client-Server Architecture:} Using Java TCP sockets (\texttt{ServerSocket}/\texttt{Socket}) for distributed booking across terminal stations.
    \item \textbf{Modern Java:} Adopting Records for immutable DTOs and Virtual Threads (Project Loom) for lightweight concurrency.
\end{enumerate}

% ==============================================================================
% REFERENCES
% ==============================================================================
\begin{thebibliography}{99}
\bibitem{bloch} 
J.~Bloch, \emph{Effective Java}, 3rd~ed.\hskip 1em plus 0.5em minus 0.4em\relax Boston, MA: Addison-Wesley, 2018.

\bibitem{schildt} 
H.~Schildt, \emph{Java: The Complete Reference}, 12th~ed.\hskip 1em plus 0.5em minus 0.4em\relax New York, NY: McGraw-Hill Education, 2021.

\bibitem{goetz} 
B.~Goetz, T.~Peierls, J.~Bloch, J.~Bowbeer, D.~Holmes, and D.~Lea, \emph{Java Concurrency in Practice}.\hskip 1em plus 0.5em minus 0.4em\relax Boston, MA: Addison-Wesley, 2006.

\bibitem{sqlitejdbc} 
Taro L. Saito, \emph{SQLite JDBC Driver Documentation}.\hskip 1em plus 0.5em minus 0.4em\relax [Online]. Available: \url{https://github.com/xerial/sqlite-jdbc}.
\end{thebibliography}

% ==============================================================================
% APPENDIX: SOURCE CODE REFERENCE
% ==============================================================================
\newpage
\appendix
\section{Source Code Reference}
\label{sec:appendix_code}

The complete implementation is contained in a single source file \texttt{Main.java} (410 lines), included with this submission. The file consolidates all classes discussed in this report. Below is a summary of the key components and their line ranges:

\begin{table}[h!]
\centering
\small
\begin{tabularx}{\textwidth}{l l p{6.5cm}}
\toprule
\textbf{Component} & \textbf{Lines} & \textbf{Description} \\
\midrule
Import statements & 1--20 & AWT, SQL, I/O, Collections, Swing imports \\
\texttt{Rentable} interface & 22--24 & Behavioral contract for rental display \\
\texttt{Vehicle} abstract class & 26--40 & Base entity with ID, model, abstract tier \\
\texttt{RentalCar} class & 42--73 & Fleet asset with category, rate, fuel class \\
\texttt{LuxurySedan} class & 75--85 & Executive specialization with chauffeur flag \\
\texttt{VehicleNotFoundException} & 87--91 & Custom checked exception \\
\texttt{FleetMaintenanceThread} & 93--97 & Thread subclass demonstration \\
\texttt{DailyBillingTask} & 99--103 & Runnable interface demonstration \\
\texttt{FleetDispatchSync} & 105--118 & Monitor sync with wait/notify \\
\texttt{Pair<K,V>} & 120--127 & Generic container class \\
\texttt{Main} class & 129--410 & JDBC CRUD, collections, main(), Swing GUI \\
\bottomrule
\end{tabularx}
\caption{Source Code Component Index for \texttt{Main.java}.}
\label{tab:code_index}
\end{table}

Key implementation excerpts have been presented inline throughout Sections~\ref{sec:structure}--\ref{sec:logic} of this report. The full source file is available as a separate submission artifact alongside this document.

\end{document}
